% ThreatLens — paper draft
% Compiles with: pdflatex main && bibtex main && pdflatex main && pdflatex main
% Uses only standard packages (no acmart/IEEEtran required).
% To switch to ACM format for CAMLIS/AISec: replace the documentclass block
% with \documentclass[sigconf]{acmart} and delete the \author block below.
%
% AUTHOR ACTION ITEMS (search for "TODO"):
%   - fill in surname / affiliation
%   - complete author fields in references.bib (entries marked TODO-AUTHORS)
%   - decide whether to include the independent-annotator study (Sec. 8)

\documentclass[10pt,twocolumn]{article}

\usepackage[margin=0.85in]{geometry}
\usepackage[T1]{fontenc}
\usepackage[utf8]{inputenc}
\usepackage{times}
\usepackage{booktabs}
\usepackage{amsmath}
\usepackage{url}
\usepackage{xcolor}
\usepackage{listings}
\usepackage[hidelinks]{hyperref}
\usepackage{balance}

\lstset{basicstyle=\ttfamily\scriptsize, breaklines=true, frame=single,
        columns=fullflexible, xleftmargin=2pt, framexleftmargin=2pt}

\newcommand{\tool}{\textsc{ThreatLens}\xspace}
\usepackage{xspace}

\title{\bf ThreatLens: Claim-Level Grounding Verification for
LLM-Extracted Cyber Threat Intelligence}

\author{
Shiddarth Dey Tusar
\\ \small Charles Sturt University, NSW, Australia
\\ \small \texttt{tusardey77@gmail.com}
}
\date{}

\begin{document}
\maketitle

\begin{abstract}
Large language models are increasingly proposed for automating cyber threat
intelligence (CTI) extraction, but existing benchmarks evaluate them with
precision-style accuracy metrics against static gold labels. We argue this
measurement approach is unsafe for operational deployment, and we present
\tool, an open-source pipeline that performs deterministic, claim-level
grounding verification of LLM-extracted CTI against three authorities: the
National Vulnerability Database, MITRE ATT\&CK STIX data, and the source
document itself. Every atomic claim (CVE, indicator of compromise, ATT\&CK
technique, threat actor, malware family) receives one of five verdicts under a
taxonomy that separates \emph{fabricated} entities from entities that are real
but \emph{ungrounded} in the source, \emph{mislabeled} metadata, and claims
that are simply \emph{unverifiable} because no authority can confirm them.

We evaluate five open-weight models on 94 CISA advisories published after the
models' training cutoffs, plus two indicator-list CTI events. Four findings
emerge. (i)~Hallucination is task-dependent rather than model-intrinsic: Nemotron-3-Super 120B
scores $0.00\%$ [95\% CI: $0.00\%, 0.52\%$] on indicator-copying tasks and $6.39\%$ [$4.65\%, 8.73\%$] on narrative
advisories requiring inference. (ii)~Precision metrics conceal
\emph{silent failure}: Nemotron-3-nano 30B returned well-formed but empty extractions on
$100.0\%$ [$72.2\%, 100.0\%$] of its successful runs over CVE-bearing advisories, which a
precision-only benchmark would report as flawless. Even Nemotron-3-Super 120B silently
returns nothing on $11.2\%$ [$7.0\%, 17.4\%$] of CVE-bearing advisories. (iii)~Two reproducible
failure modes appear at temperature~0: parametric-memory CVE imports and
systematic ATT\&CK sub-technique confabulation. (iv)~Following NIST's April 2026
decision to abandon universal NVD enrichment, existence checking against
vulnerability registries fails precisely on the freshest advisories: 9 of 11
CVEs cited verbatim in same-week ICS advisories were absent from both NVD and
CVE.org. We further report the verifier's own audit trail: five classes of
verifier false positive discovered by manual adjudication and fixed under
regression test, which we argue should be a standard reporting requirement for
automated CTI evaluation. Code, corpus manifests, and per-claim verdicts are
released.
\end{abstract}

% ============================================================================
\section{Introduction}

Security operations centres face more threat reporting than analysts can read.
Large language models (LLMs) are an appealing remedy: they can read an advisory
and emit structured indicators, techniques, and actor attributions in seconds.
A growing body of work benchmarks this capability
\cite{alam2024ctibench,alam2025athenabench,cyberthreateval2026}, and a
complementary line of work documents its unreliability
\cite{mezzi2025unreliable,uncovering2025vulnerabilities}.

Our concern is with how that unreliability is \emph{measured}. Existing CTI
benchmarks score model outputs against static gold labels and report
accuracy-style aggregates. Three problems follow. First, gold labels age: the
threat landscape moves faster than benchmark curation, and static datasets
leak into training corpora. Second, an accuracy score over a fixed answer key
says nothing about whether a specific claim in a specific report is safe to act
on, which is the question an analyst actually asks. Third, and most seriously,
precision-style metrics are blind to omission: a model that extracts nothing
cannot be wrong.

We take a different position. Rather than scoring models against an answer key,
we verify each extracted claim against the authorities that operational
security already trusts, at extraction time, deterministically, and without a
human in the loop. This reframes evaluation as \emph{grounding} rather than
\emph{accuracy}, and it yields a tool that is useful in a pipeline, not only in
a leaderboard.

\paragraph{Contributions.}
\begin{itemize}\itemsep2pt
\item \tool, an open-source claim-level CTI verifier that cross-checks LLM
      output against NVD \cite{nvdapi}, CVE.org \cite{cveorg}, MITRE ATT\&CK
      STIX \cite{attackstix}, and source-document provenance
      (Section~\ref{sec:design}).
\item A five-verdict taxonomy that distinguishes fabricated entities from
      real-but-ungrounded ones---a distinction absent from prior CTI
      evaluation and essential for diagnosing parametric-memory leakage.
\item An evaluation on 94 post-cutoff CISA advisories showing that
      hallucination rates are task-dependent, that silent failure is prevalent
      and invisible to precision metrics, and that two specific failure modes
      reproduce deterministically (Section~\ref{sec:results}).
\item Quantification of vulnerability-registry lag as a verification failure
      mode, in the wake of NIST's 2026 curtailment of NVD enrichment
      \cite{nist2026nvd,helpnet2026nvd}.
\item An explicit verifier audit trail: five false-positive classes found by
      adjudicating flagged claims, each fixed and regression-tested. We argue
      this should be reported as standard practice
      (Section~\ref{sec:audit}).
\end{itemize}

% ============================================================================
\section{Background and Related Work}

\paragraph{CTI benchmarks.}
CTIBench \cite{alam2024ctibench} evaluates LLMs across CTI knowledge questions,
CVE-to-CWE mapping, severity prediction, ATT\&CK technique extraction and
actor attribution. AthenaBench \cite{alam2025athenabench} extends it with a
dynamic dataset pipeline, deduplication and a risk-mitigation task, finding
open models substantially behind proprietary ones on reasoning-heavy tasks.
Both are answer-key benchmarks: they measure whether the model matches a label,
not whether a claim is grounded in the document that produced it.

\paragraph{Reliability studies.}
Mezzi et al.\ \cite{mezzi2025unreliable} evaluate consistency and confidence
over 350 threat reports and find LLMs inconsistent and overconfident on
real-size inputs. \cite{uncovering2025vulnerabilities} categorises CTI failure
modes---spurious metadata correlations, contradictory sources, and constrained
generalisation to emerging threats---using a human-in-the-loop framework. These
works diagnose; neither ships an automated verifier. Feng and Sakurai
\cite{feng2025nids} survey LLM integration in intrusion detection and its
attendant risks.

\paragraph{Closest work.}
CyberThreat-Eval \cite{cyberthreateval2026} is an expert-annotated benchmark
drawn from a production CTI workflow, with analyst-centric metrics; its agent
integrates external ground truth and expert feedback. It observes, as we do,
that LLMs ``struggle to distinguish correct from incorrect information'' and
that external ground truth for technical details is often absent. It differs
from \tool in three ways: it is built around a proprietary corporate workflow
with expert annotators, the verification layer is not released as a standalone
reproducible artifact, and it evaluates frontier proprietary models rather than
the open models a resource-constrained SOC can actually run.

\paragraph{Grounding and hallucination detection.}
General-purpose grounding frameworks such as GSAR \cite{gsar2026} partition
claims into grounded/ungrounded typologies for multi-agent LLM systems, and
domain-grounded retrieval pipelines perform atomic claim verification
\cite{tieredretrieval2026}; neither is specialised to CTI authorities. KGV
\cite{kgv2024} assesses CTI credibility with knowledge graphs, targeting source
credibility rather than claim-level grounding of extractions. To our knowledge
no open-source tool performs deterministic claim-level verification of
LLM-extracted CTI against vulnerability and adversary-technique registries.

% ============================================================================
\section{ThreatLens Design}
\label{sec:design}

\subsection{Claim model}
An extraction is decomposed into atomic \emph{claims}. Each claim carries a
type (CVE, IoC-IP, IoC-domain, IoC-URL, IoC-hash, IoC-email, ATT\&CK technique,
threat actor, malware), a value, an optional label (product name, technique
name), and an evidence span the model is instructed to quote verbatim from the
source.

\subsection{Verdict taxonomy}
Each claim receives exactly one verdict:

\begin{description}\itemsep2pt
\item[\textsc{Verified}] the entity exists in an authority, is grounded in the
  source document, and its metadata is consistent.
\item[\textsc{Fabricated}] the entity does not exist in any authority and is
  not present in the source (e.g.\ a malformed indicator, an invalid ATT\&CK
  identifier, an unregistered and unquoted CVE).
\item[\textsc{Ungrounded}] the entity is real but the source document does not
  support it. This isolates \emph{parametric-memory import}: the model supplies
  a fact from training rather than from the document.
\item[\textsc{Mislabeled}] the entity is real and grounded but the attached
  metadata is wrong (e.g.\ a valid technique identifier paired with the name of
  a different technique).
\item[\textsc{Unverifiable}] no authority can adjudicate the claim---the API is
  unreachable, or the entity is source-attested but absent from all registries.
  Verdicts are never guessed; this is a fail-safe state, not a failure.
\end{description}

The \textsc{Ungrounded} / \textsc{Fabricated} split is the taxonomy's key
contribution. A real CVE that the source never mentions is not an invented
fact, but it is equally unsafe in an automated pipeline: it is an
unattributable assertion. Prior CTI evaluation collapses both into
``incorrect''.

\subsection{Authorities}
\textbf{Vulnerabilities.} CVEs are checked against NVD API~2.0 \cite{nvdapi}
with an on-disk cache and rate limiting. Because NVD enrichment is now
incomplete by policy \cite{nist2026nvd}, a negative NVD result triggers a
secondary lookup against CVE.org CNA records \cite{cveorg}. Only a claim absent
from both \emph{and} absent from the source text is \textsc{Fabricated}.

\textbf{Adversary behaviour.} Technique identifiers, intrusion-set aliases and
malware/tool names are checked against the MITRE ATT\&CK Enterprise STIX~2.1
bundle \cite{attackstix,oasis2021stix,strom2018attack} (v18.1: 835 techniques,
580 actor aliases, 1{,}062 malware and tool names).

\textbf{Provenance.} Every claim type is additionally checked for literal
presence in the source document after defanging normalisation
(\texttt{hxxp}$\rightarrow$\texttt{http}, \texttt{[.]}$\rightarrow$\texttt{.}).
Indicators must appear verbatim: there is no legitimate mechanism by which a
model can know a C2 address or file hash that the report does not state.
Technique identifiers are treated differently, since they are analyst
inferences rather than quotations; for these the model's evidence span must
appear in the source.

\subsection{Pipeline}
Ingestion normalises CISA advisories \cite{cisaadvisories} and MISP OSINT
events \cite{wagner2016misp,circlfeed} into a common schema with publication
dates. A provider-agnostic client supports local Ollama \cite{ollama} and any
OpenAI-compatible endpoint \cite{openrouter}. Extraction uses a fixed
JSON-schema prompt at temperature~0; the pipeline never filters model output,
so raw behaviour is measured. Verification is deterministic, and results
(including every per-claim check) are appended to JSONL with resume support.
A separate re-verification tool recomputes all verdicts under revised rules
without re-invoking any model, which makes verifier refinement cheap and
auditable.

% ============================================================================
\section{Experimental Setup}
\label{sec:setup}

\paragraph{Corpora.}
The primary corpus is 94 unique CISA advisories \cite{cisaadvisories} collected
between 2026-06-12 and 2026-07-25 (ICS advisories, joint alerts, and KEV
catalogue additions), mean length 13{,}104 characters, containing 285 distinct
CVE mentions. All were published after the evaluated models' training cutoffs,
which removes memorisation as an explanation for correct extraction. A
secondary corpus of two TLP:WHITE CIRCL MISP OSINT events \cite{circlfeed}
(Turla infrastructure; Operation Dust Storm; $\approx$235 indicators, mean
3{,}754 characters) provides an indicator-copying contrast condition.

\paragraph{Models.}
We attempted twelve open-weight endpoints via OpenRouter's API \cite{openrouter}; five produced usable output
(Table~\ref{tab:main}): \texttt{nvidia/nemotron-3-super-120b-a12b:free}, \texttt{openai/gpt-oss-20b:free}, \texttt{google/gemma-4-31b-it:free}, \texttt{nvidia/nemotron-nano-9b-v2:free}, and \texttt{nvidia/nemotron-3-nano-30b-a3b:free} (evaluated between June 12 and July 25, 2026). Seven endpoints produced no usable extractions due to persistent upstream HTTP 429 rate-limiting or endpoint withdrawal---a constraint we report rather than hide (Section~\ref{sec:limits}).

\paragraph{Protocol.}
Temperature~0, fixed prompt, up to three runs per report per model. Claims are
never filtered before verification. A deliberately corrupted extraction with
injected errors of every taxonomy class is carried through the pipeline as a
negative control on every re-verification.

% ============================================================================
\section{Results}
\label{sec:results}

\begin{table}[t]
\centering\small
\caption{Grounding verdicts by model and corpus. \textsc{Cisa} = narrative
advisories (inference required); \textsc{Misp} = indicator lists (copying).
Hallucination rate $=$ (fabricated $+$ ungrounded $+$ mislabeled) $/$
verifiable claims, reported with Wilson 95\% CIs.}
\label{tab:main}
\begin{tabular}{llrrrr}
\toprule
Model & Corpus & Runs & Claims & Halluc. Rate & 95\% CI \\
\midrule
Nemotron-3-Super 120B & \textsc{misp} & 6   & 739 & \bf 0.00\% & [0.00\%, 0.52\%] \\
Nemotron-3-Super 120B & \textsc{cisa} & 143 & 590 & \bf 6.39\% & [4.65\%, 8.73\%] \\
gpt-oss-20B           & \textsc{misp} & 1   & 42  & \bf 0.00\% & [0.00\%, 8.38\%] \\
gpt-oss-20B           & \textsc{cisa} & 39  & 113 & \bf 23.71\% & [16.35\%, 33.07\%] \\
Gemma-4 31B           & \textsc{misp} & 1   & 39  & 0.00\% & [0.00\%, 8.97\%] \\
Nemotron-nano 9B      & \textsc{cisa} & 5   & 7   & 0.00\% & [0.00\%, 35.43\%] \\
Nemotron-3-nano 30B   & \textsc{cisa} & 10  & 0   & --- & --- \\
\midrule
Negative control      & \textsc{misp} & 1   & 8   & 87.50\% & [52.91\%, 97.76\%] \\
\bottomrule
\end{tabular}
\end{table}

\subsection{Hallucination is task-dependent}
Table~\ref{tab:main} shows the same model, prompt and verifier producing
$0.00\%$ [95\% CI: $0.00\%, 0.52\%$] hallucination on indicator-list events and $6.39\%$ [$4.65\%, 8.73\%$] on narrative
advisories for Nemotron-3-Super 120B, and $0.00\%$ [$0.00\%, 8.38\%$] versus $23.71\%$ [$16.35\%, 33.07\%$] for gpt-oss-20B. The
distinction is not report length but task type: indicator lists require
copying, whereas advisories require inferring techniques and mapping products,
and inference is where grounding fails. A benchmark composed mainly of
extraction-from-list tasks will therefore report reassuring numbers that do not
transfer to the narrative reports analysts actually triage.

\subsection{Silent failure is invisible to precision metrics}
\label{sec:silent}

\begin{table}[t]
\centering\footnotesize
\caption{Yield analysis over CVE-bearing advisories with Wilson 95\% CIs. \emph{Silent failure} =
well-formed but empty extraction from a report containing CVEs;
\emph{abstention} = empty extraction from a report with none.}
\label{tab:silent}
\begin{tabular}{lrrrr}
\toprule
Model & Runs & Yield [95\% CI] & Silent [95\% CI] & Abstain [95\% CI] \\
\midrule
Nemotron-3-Super 120B & 143 & 86.7\% [80.2\%, 91.3\%] & \bf 11.2\% [7.0\%, 17.4\%] & 2.1\% [0.7\%, 6.0\%] \\
gpt-oss-20B           & 39  & 97.4\% [86.8\%, 99.5\%] & 2.6\% [0.5\%, 13.2\%]     & 0.0\% [0.0\%, 9.0\%] \\
Nemotron-nano 9B      & 5   & 80.0\% [37.6\%, 96.4\%] & 0.0\% [0.0\%, 43.4\%]    & 20.0\% [3.6\%, 62.4\%] \\
Nemotron-3-nano 30B   & 10  & 0.0\% [0.0\%, 27.8\%]   & \bf 100.0\% [72.2\%, 100.0\%] & 0.0\% [0.0\%, 27.8\%] \\
\bottomrule
\end{tabular}
\end{table}

Nemotron-3-nano 30B returned syntactically valid JSON containing empty arrays
on \emph{every} successful run ($100.0\%$ [95\% CI: $72.2\%, 100.0\%$]). Its hallucination rate is undefined because it
made no claims; a precision-only benchmark would report it as making no errors.
Table~\ref{tab:silent} separates this from legitimate abstention by asking
whether the source document contained CVE identifiers at all. All ten of Nemotron-3-nano 30B's empty runs were on advisories that did contain CVEs---up to seven
in one case---so they are unambiguous failures, not appropriate silence.

The result generalises beyond the smallest models: the best-performing model in our study, Nemotron-3-Super 120B, silently returns nothing on $11.2\%$ [$7.0\%, 17.4\%$] of CVE-bearing advisories
while reporting an excellent $6.39\%$ [$4.65\%, 8.73\%$] hallucination rate. In an automated
defence pipeline a silently dropped advisory is arguably more dangerous than a
visible error, because nothing signals that review is needed. We therefore
argue that \emph{yield} and \emph{silent-failure rate} must be reported
alongside any precision-style CTI metric.

\subsection{Two reproducible failure modes}

\paragraph{Parametric-memory CVE import.}
On KEV-catalogue alerts whose page text announces additions without listing
identifiers, models supplied CVE identifiers that are real, correctly
formatted, and absent from the source (e.g.\ CVE-2026-10520, CVE-2026-42271,
CVE-2026-50751; 5 of 5 runs). These are \textsc{Ungrounded}: the model answered
from training memory rather than from the document. Under a conventional
answer-key benchmark such a claim may well be scored \emph{correct}, since the
CVE is genuinely associated with the alert. Under operational conditions it is
an unattributable assertion whose provenance cannot be traced to the input.

\paragraph{ATT\&CK sub-technique confabulation.}
Models pair valid technique identifiers with the names of neighbouring
sub-techniques: T1071.001 (\emph{Web Protocols}) labelled ``Application Layer
Protocol: MQTT'' (correctly T1071.005), and T1087.002 (\emph{Domain Account})
labelled ``Cloud Account Discovery'' (correctly T1087.004). These recurred
identically across all runs at temperature~0, so they are systematic properties
of the model rather than sampling noise. Because both components are
individually plausible, downstream consumers keyed on the identifier would
silently ingest a wrong technique.

\subsection{Registry lag defeats existence checking}
NIST curtailed universal NVD enrichment in April 2026, prioritising only
KEV-listed, federally used, and critical software \cite{nist2026nvd,
helpnet2026nvd}. We observe the operational consequence directly: across three
same-week ICS advisories, 9 of 11 CVEs quoted \emph{verbatim} in the advisory
text were present in neither NVD nor CVE.org, while the remaining two resolved
normally. A naive verifier flags all nine as fabrications. The correct verdict
is \textsc{Unverifiable}: the model was faithful, and the registries were
merely behind. Over the full study, 27 claims from Nemotron-3-Super 120B were
quarantined this way. Existence checking against vulnerability registries thus
degrades exactly on the freshest---and most operationally urgent---advisories,
which we believe is an under-appreciated constraint on automated CTI validation
generally.

% ============================================================================
\section{The Verifier Audit Trail}
\label{sec:audit}

A verifier that is itself wrong produces confident nonsense. We therefore
adjudicated every flagged claim by hand and treated verifier false positives as
first-class findings. Five classes emerged, each fixed and pinned by regression
tests (23 tests at time of writing):

\begin{enumerate}\itemsep2pt
\item \textbf{Indicator type confusion.} A model extracted an X.509 SHA-1
      fingerprint line into an \texttt{ip} field. The artifact is real and
      grounded; only the type is wrong. Rule: grounded-but-malformed indicators
      are \textsc{Mislabeled}, not \textsc{Fabricated}.
\item \textbf{Registry lag} (Section~\ref{sec:results}): grounded but
      unregistered CVEs are \textsc{Unverifiable}.
\item \textbf{Parent-prefixed technique labels.} ``Unsecured Credentials:
      Private Keys'' is correct usage for T1552.004 but failed naive string
      similarity against the official sub-technique name.
\item \textbf{Product-name comparison.} Comparing a short product label against
      a long CVE description with sequence similarity systematically
      under-scores correct labels; token-overlap comparison fixed this.
\item \textbf{Supply-chain product mismatch.} ICS advisories name the
      \emph{affected product} (e.g.\ Hitachi Energy RTU500) while NVD describes
      the \emph{vulnerable component}. A label grounded in the advisory is
      faithful extraction; we mark these \textsc{Verified} and flag them, which
      incidentally yields a dataset of supply-chain vulnerability cases.
\end{enumerate}

Before these fixes the measured hallucination rate for Nemotron-3-Super 120B was
$17.5\%$; after adjudication it is $6.39\%$ [$4.65\%, 8.73\%$] on the same data. The difference is
entirely verifier error. We report this trajectory deliberately: any automated
CTI evaluation that has not published such an audit should be assumed to
contain false positives of similar magnitude.

% ============================================================================
\section{Discussion}

Three implications follow for practitioners. First, precision is not safety:
deployment decisions require yield and silent-failure rates, and a model that
abstains on hard inputs is not thereby safe. Second, verification must be
multi-authority and fail-safe, because single-source existence checks now break
on recent advisories by policy rather than by accident. Third, the
grounded/ungrounded distinction should be enforced in pipelines: a real CVE
that the source never mentions must be surfaced for review, not silently
promoted to an action.

For researchers, our results suggest that CTI benchmark scores are not
comparable across task types, and that reporting a single hallucination rate
per model obscures the operative variable, which is what the task demands of
the model rather than what the model is.

% ============================================================================
\section{Limitations}
\label{sec:limits}

Our study is preliminary in scale and we state its bounds plainly.
The corpus is 94 advisories dominated by ICS bulletins from a single publisher;
adversary-narrative reports from vendor research would exercise actor and
technique extraction more heavily. Sample sizes are unequal and small for three
of five models (5--39 successful runs), resulting in wide Wilson 95\% confidence intervals (e.g., $0.00\%$ [$0.00\%, 35.43\%$] for Nemotron-nano 9B and $23.71\%$ [$16.35\%, 33.07\%$] for gpt-oss-20B); only Nemotron-3-Super 120B has sufficient sample size ($n=563$ verifiable claims, 95\% CI [$4.65\%, 8.73\%$]) for tight point estimation. Model availability, not compute, was the binding constraint: across the
study we recorded 305 dropped connections and roughly 120 rate-limit rejections
on free endpoints, and 7 of 12 attempted models produced no usable output at
all. Reproducibility of free-tier evaluation is therefore poor, and we
recommend local inference \cite{ollama} or paid endpoints for future work.

Verification covers only claim types with machine-readable authorities.
Narrative assertions---attribution reasoning, impact assessment, campaign
linkage---remain outside scope, and our hallucination rates should not be read
as covering them. Adjudication of flagged claims was performed by one annotator
(the author); an independent-annotator study with inter-rater agreement is
required before the verifier's own precision can be stated numerically, and we
regard this as the most important piece of remaining work. Finally, runs were
conducted at temperature~0 with a single prompt; prompt sensitivity and
retrieval-augmented variants are unexplored.

% ============================================================================
\section{Availability}

\tool is released with the verifier, ingestion modules for CISA and MISP
sources, the corpus manifest with publication dates, per-claim verdict records,
the regression test suite, and the re-verification tool that reproduces every
number in this paper from stored extractions without re-invoking any model.
% TODO: insert repository URL on release
\begin{center}\small\url{https://github.com/ShiddarthDey/threatlense.git}\end{center}

\balance
\bibliographystyle{plain}
\bibliography{references}

\end{document}
